\documentclass[a4paper,10pt,fleqn]{article}
\usepackage{fuzz}
\begin{document}
\title{Vox voice selection}
\maketitle

A roster belongs to a target (a provider and its optional model).
The panel may lag behind the file. The saved voice must belong to the
saved target's roster; a hand-edited invalid target is represented by
an empty roster and no confirmed voice. Picker indices name entries in
the shown roster, never in a different target's roster.

\begin{zed}
  [TARGET, VOICE]
\end{zed}

\begin{schema}{VoiceSelection}
  available : TARGET \fun \power VOICE \\
  saved, shown : TARGET \\
  voice, roster : \power VOICE
\where
  \# voice \leq 1 \\
  voice \subseteq available(saved) \\
  roster \subseteq available(shown)
\end{schema}

\begin{schema}{Init}
  VoiceSelection
\where
  shown = saved \\
  voice = \emptyset \\
  roster = \emptyset
\end{schema}

\begin{schema}{SwitchTarget}
  \Delta VoiceSelection \\
  target? : TARGET \\
  default? : \power VOICE
\where
  available' = available \\
  saved' = target? \\
  voice' = default? \\
  shown' = shown \\
  roster' = roster
\end{schema}

A panel model click also installs the fetched roster, as in a successful
refresh below. A daemon rejection can show the configured target without
offering any voices; the model picker remains usable. Roster order is left
abstract: a voice pick enters the model after its index has been decoded.

\begin{schema}{PanelSwitchTarget}
  \Delta VoiceSelection \\
  target?, expected? : TARGET \\
  fetched?, default? : \power VOICE
\where
  expected? = saved \\
  default? \subseteq fetched? \\
  available' = available \\
  shown' = target? \\
  saved' = target? \\
  roster' = fetched? \\
  voice' = default?
\end{schema}

\begin{schema}{RefreshRoster}
  \Delta VoiceSelection \\
  target? : TARGET \\
  fetched? : \power VOICE
\where
  target? = saved \\
  available' = available \\
  saved' = saved \\
  voice' = voice \\
  shown' = target? \\
  roster' = fetched?
\end{schema}

\begin{schema}{ChooseVoice}
  \Delta VoiceSelection \\
  picked? : VOICE
\where
  shown = saved \\
  picked? \in roster \\
  voice' = \{picked?\} \\
  available' = available \\
  saved' = saved \\
  shown' = shown \\
  roster' = roster
\end{schema}

In a stale click (\texttt{shown} differs from \texttt{saved}), the choice is rejected
without a write; a refresh restores the current roster. A stale refresh
similarly leaves the confirmed settings unchanged.

\begin{schema}{RejectStaleVoice}
  \Xi VoiceSelection
\where
  shown \neq saved
\end{schema}

\begin{schema}{RejectStaleRefresh}
  \Xi VoiceSelection \\
  target? : TARGET
\where
  target? \neq saved
\end{schema}

\end{document}
